\section{Quiz}
\label{sec:quiz}

To reinforce the concepts discussed during the seminar, the following multiple-choice questions have been designed. They cover the fundamental mechanics of the Classic Visitor pattern, as well as the architectural traits of its modern variants.

\begin{enumerate}
    \item \textbf{What type of dispatch is naturally applied by C++, and what type of dispatch is simulated by the Visitor pattern?}
    \begin{enumerate}[label=\textbf{\Alph*.}]
        \item Double-dispatch, single-dispatch
        \item Both single-dispatch
        \item Both double-dispatch
        \item Single-dispatch, double-dispatch
    \end{enumerate}
    \textit{Answer: D}

    \item \textbf{When is the BEST time to use the Visitor pattern?}
    \begin{enumerate}[label=\textbf{\Alph*.}]
        \item When the classes defining the object structure change frequently.
        \item When the object structure is completely stable, but you often need to define new operations over it.
        \item When you need to hide the creation logic of an object.
        \item When you want to pass a request along a chain of objects until it is handled.
    \end{enumerate}
    \textit{Answer: B}

    \item \textbf{What is one of the main negative consequences (drawbacks) of using the Visitor pattern?}
    \begin{enumerate}[label=\textbf{\Alph*.}]
        \item It makes adding new operations incredibly difficult.
        \item It can break encapsulation by forcing elements to expose their internal state via public methods.
        \item It prevents visitors from accumulating state during traversal.
        \item It forces all elements to share a single, strict parent class.
    \end{enumerate}
    \textit{Answer: B}
\end{enumerate}


\begin{enumerate}
    \setcounter{enumi}{3} % Bắt đầu đếm tiếp từ câu 4
    \item \textbf{How does the modern \texttt{std::variant} approach eliminate the need for an abstract base class and the intrusive \texttt{accept} method?}
    \begin{enumerate}[label=\textbf{\Alph*.}]
        \item By using \texttt{dynamic\_cast} to resolve types at runtime recursively.
        \item By relying on compile-time type-based visitation over a closed set of alternatives and invoking an overloaded \texttt{operator()} inside a Functor.
        \item By implementing the double-dispatch mechanism exclusively through \texttt{virtual} functions.
        \item By forcing all elements to inherit from a common \texttt{std::variant} interface.
    \end{enumerate}
    \textit{Answer: B}
    
    \item \textbf{What is the primary architectural technique utilized by the Acyclic Visitor pattern to break the circular dependency between elements and visitors?}
    \begin{enumerate}[label=\textbf{\Alph*.}]
        \item Utilizing a single degenerate (empty) base interface combined with runtime cross-casting (\texttt{dynamic\_cast}).
        \item Extracting all \texttt{visit} methods into a standalone \texttt{std::visit} function.
        \item Using multiple inheritance for the Element classes instead of the Visitor classes.
        \item Storing elements contiguously in memory using \texttt{std::vector} to avoid pointer dependencies.
    \end{enumerate}
    \textit{Answer: A}

    \item \textbf{According to the benchmark, what primarily contributes to the observed runtime advantage of the \texttt{std::variant} Visitor over the Classic and Acyclic Visitors?}
    \begin{enumerate}[label=\textbf{\Alph*.}]
        \item It utilizes RTTI lookups (\texttt{typeid}), which are heavily optimized by modern CPUs.
        \item It executes fewer instructions by avoiding \texttt{dynamic\_cast} while continuing to use heap-allocated pointers.
        \item It runs in a multi-threaded context automatically via standard C++17 library optimizations.
        \item It stores objects by value, maximizing cache locality (cache hits), and entirely avoids virtual table (vtable) indirection.
    \end{enumerate}
    \textit{Answer: D}

    \item \textbf{If new element types are expected to be added frequently, which approach is generally more suitable?}
    \begin{enumerate}[label=\textbf{\Alph*.}]
        \item Classic Visitor, because it strictly adheres to the Open/Closed Principle for both operations and types.
        \item \texttt{std::variant}, because it easily supports adding new types without modifying the variant alias or recompiling.
        \item Acyclic Visitor, because it allows introducing new element-specific visitor interfaces without modifying a monolithic Visitor interface.
        \item None of them, as the Visitor pattern is fundamentally flawed when dealing with new operations.
    \end{enumerate}
    \textit{Answer: C}
\end{enumerate}
